\section{Interventional Framework}

For each base task $b$, \benchmark constructs a clean instance $x_b^0$ and intervention instances $x_b^j$ for intervention families $j$. Each intervention is intended to change one targeted factor while keeping the user's high-level goal fixed. We use \emph{interventional} or \emph{paired-perturbation} language for this controlled benchmark design; it is inspired by causal and counterfactual evaluation \citep{holland1986statisticscausalinference,pearl2009causality} but does not constitute observational causal identification or deployment-wide effect estimation. The intervention metadata records the family, target factor, non-target factors expected to remain stable, expected robust behavior, severity, patch details, and whether the final answer is expected to change.

The current intervention families are tool removal, tool failure, tool corruption, irrelevant tools, memory corruption, observation conflict, ambiguous instruction, long-horizon dependency, premature success signal, and distractor evidence. These families target the skill components described in the claim ledger: recovery under failures (\claimref{C2}, \claimref{C5}), contradiction handling (\claimref{C3}, \claimref{C8}), memory verification (\claimref{C2}, \claimref{C3}), tool overuse (\claimref{C7}), and stopping behavior (\claimref{C8}).

The causal interpretation depends on intervention validity, which \benchmark audits rather than assumes. A tool-failure condition should primarily affect tool reliability, not silently alter the user's objective or the correct answer unless explicitly specified. A memory-corruption condition should make memory less trustworthy while still allowing verification. An observation-conflict condition should create a conflict that can be resolved or acknowledged, not an impossible task unless impossibility is part of the expected robust behavior. The benchmark audit flags interventions that change multiple patch groups, remove required tools outside the target family, lack expected-behavior or scoring metadata, leave ground-truth changes undocumented, or fall outside family-specific severity ranges. Expert or human audit remains necessary, and this remaining risk is tracked by \claimref{C10}.

\paragraph{Automated isolation audit.}
The repository includes an intervention-isolation audit (\texttt{scripts/audit\_intervention\_isolation.py}) that checks whether paired clean/intervention instances change only the intended patch groups. Audit reports are versioned under \texttt{audits/}. Automated pass is necessary but not sufficient for \claimref{C10}; expert or human review of a stratified sample remains required before claiming that interventions isolate skill components.
